\chapter{Machine-learning methods}
\label{app:methods}

This appendix defines the standard machine-learning components used in
Chapter~\ref{ch:arch}, gives each one's formula, and states exactly how it is
configured in this work. Symbols: $\theta$ are the model parameters, $g_t=\nabla_\theta
\mathcal{L}$ the gradient at step $t$, and $\eta$ the learning rate.

\section{Adam optimiser}
Adam \citep{kingma2015} adapts a per-parameter step size from running estimates of
the first and second moments of the gradient:
\begin{align}
  m_t &= \beta_1 m_{t-1} + (1-\beta_1)\,g_t, &
  v_t &= \beta_2 v_{t-1} + (1-\beta_2)\,g_t^2, \nonumber\\
  \hat{m}_t &= \frac{m_t}{1-\beta_1^{\,t}}, &
  \hat{v}_t &= \frac{v_t}{1-\beta_2^{\,t}}, \qquad
  \theta_t = \theta_{t-1} - \eta\,\frac{\hat{m}_t}{\sqrt{\hat{v}_t}+\epsilon}.
\end{align}
\textbf{Implementation.} \texttt{torch.optim.Adam} over all $3.3$\,M parameters,
$\eta=10^{-3}$, $(\beta_1,\beta_2)=(0.9,0.999)$, $\epsilon=10^{-8}$; one update per
mini-batch of $8$ windows.

\section{Cosine learning-rate decay}
The learning rate is annealed from its initial value to zero along a half cosine
over the training budget of $T$ epochs \citep{loshchilov2017}:
\begin{equation}
  \eta_t = \tfrac{1}{2}\,\eta_0\left(1 + \cos\frac{\pi t}{T}\right).
\end{equation}
\textbf{Implementation.} \texttt{CosineAnnealingLR} with $\eta_0=10^{-3}$ and
$T=130$ epochs, stepped once per epoch; large steps early, near-zero steps as
training ends.

\section{Group normalisation}
Group normalisation \citep{wu2018} splits the $C$ channels of a feature map into
$G$ groups and standardises each group over its channels and spatial extent
(independently per sample), then applies a learned affine map:
\begin{equation}
  \hat{x} = \frac{x-\mu_g}{\sqrt{\sigma_g^2+\epsilon}},\qquad y=\gamma\,\hat{x}+\beta,
\end{equation}
where $\mu_g,\sigma_g^2$ are the mean and variance over group $g$.
\textbf{Implementation.} \texttt{GroupNorm} with $G=8$ groups over the $C=96$
hidden channels ($12$ channels per group), $\epsilon=10^{-5}$, applied between
stacked ConvLSTM layers. It keeps \emph{no running statistics}, so it behaves
identically inside the unrolled autoregressive decode loop---the reason it is
preferred here over batch normalisation.

\section{Huber (smooth-$L_1$) loss}
The training criterion on the residual $r=\hat{x}-x$ is quadratic near zero and
linear in the tails, with elbow $\beta$ \citep{huber1964}:
\begin{equation}
  \ell_\beta(r)=
  \begin{cases}
    \tfrac{1}{2}\,r^2/\beta, & |r|<\beta,\\[2pt]
    |r|-\tfrac{1}{2}\beta, & |r|\ge\beta.
  \end{cases}
\end{equation}
\textbf{Implementation.} \texttt{SmoothL1Loss} with $\beta=1$, averaged over all
pixels, the four horizon frames, and the batch. The quadratic core gives stable
gradients for small residuals; the linear tail stops the rare high-amplitude
frames from dominating the update.

\section{Gradient-norm clipping}
To curb the occasional exploding gradient of a recurrent model \citep{pascanu2013},
the full gradient vector is rescaled whenever its $L_2$ norm exceeds a threshold $c$:
\begin{equation}
  g \;\leftarrow\; g\cdot\min\!\left(1,\ \frac{c}{\lVert g\rVert_2}\right).
\end{equation}
\textbf{Implementation.} \texttt{clip\_grad\_norm\_} with $c=1.0$, computed over the
concatenation of all parameter gradients, applied after back-propagation and before
each optimiser step.

\section{Bootstrap confidence interval}
The skill's uncertainty is estimated non-parametrically by resampling the
evaluation windows \citep{efron1979}. From the $n=46$ held-out windows, draw $n$
with replacement, recompute skill on the resample,
\begin{equation}
  \text{skill}^\ast_b = 1 - \frac{\operatorname{mean}_{i\in S_b}\mathrm{MAE}^{\text{model}}_i}
                             {\operatorname{mean}_{i\in S_b}\mathrm{MAE}^{\text{pers}}_i},
  \qquad b=1,\dots,B,
\end{equation}
and take the $2.5$th and $97.5$th percentiles of $\{\text{skill}^\ast_b\}$ as the
$95\%$ interval.
\textbf{Implementation.} $B=20{,}000$ resamples of whole windows (so within-window
pixel correlation is respected), giving $[+6.8\%,+10.3\%]$ and
$P(\text{skill}>0)=1.000$ for the result of Table~\ref{tab:results}.
